\section{Finite affine and path algebra}
\label{sec:affine}

Everything needed from Euclidean geometry reduces to polynomial identities in two coordinates and recursion over finite lists.

\subsection{Points, translations, and half-turns}

For points or vectors \(x=(x_1,x_2)\) and \(y=(y_1,y_2)\), write
\[
 \det(x,y)=x_1y_2-x_2y_1.
\]
For a vector \(w\), let \(\trans{w}(x)=x+w\).  For a doubled center \(s\), let
\[
 \half{s}(x)=s-x.
\]
Thus \(\half{s}\) is rotation by \(180^\circ\) about \(s/2\), but no midpoint division occurs in the formal language.

\begin{lemma}[Half-turn algebra]
\label{lem:halfturnalgebra}
For all doubled centers \(p,q\) and all points \(x\),
\begin{align}
 \half{p}(\half{p}(x))&=x,\label{eq:hinv}\\
 \half{q}(\half{p}(x))&=\trans{q-p}(x),\label{eq:hh}\\
 \half{p}(\trans{w}(x))&=\trans{-w}(\half{p}(x)).\label{eq:conjugate}
\end{align}
\end{lemma}
\begin{proof}
Expand the definitions:
\[
 p-(p-x)=x,
 \qquad
 q-(p-x)=x+(q-p),
 \qquad
 p-(x+w)=(p-x)-w.
\]
These are ring identities over the integers.  Their iPRA translations are handled by the integer-difference layer and \code{ipra\_ring}.
\end{proof}

For an oriented finite path \(P=(p_0,\ldots,p_r)\), transformations act pointwise.  Its reverse is
\[
 \rev(P)=(p_r,\ldots,p_0).
\]
As an oriented one-chain we also write \(-P\) for \(\rev(P)\).  Concatenation is denoted by \(P*Q\) when the last vertex of \(P\) is the first vertex of \(Q\).

\subsection{Doubled centers forced by endpoints}

Introduce the doubled centers of the four centrosymmetric paths:
\begin{equation}
 p=B+C,
 \qquad q=C+D,
 \qquad r=E+F,
 \qquad s=F+A.
 \label{eq:centers}
\end{equation}
The translated-pair condition gives
\begin{equation}
 E=A+t,
 \qquad D=B+t.
 \label{eq:endpointtranslation}
\end{equation}
Hence
\begin{equation}
 q=p+t,
 \qquad r=s+t.
 \label{eq:centerrelations}
\end{equation}
The second period in \cref{eq:periods} is simply
\begin{equation}
 u=r-p.
 \label{eq:u-centers}
\end{equation}
Combining \cref{eq:hh,eq:centerrelations,eq:u-centers} yields
\begin{align}
 \half{r}\half{p}&=\trans{u},
 &\half{r}\half{q}&=\trans{u-t},
 &\half{r}\half{r}&=\operatorname{id},
 &\half{r}\half{s}&=\trans{t}.
 \label{eq:fourcompositions}
\end{align}
These four identities are the whole group-theoretic content of the construction.

\subsection{Oriented path identities}

The Conway hypotheses say
\[
 \half{p}(b)=-b,
 \quad \half{q}(c)=-c,
 \quad \half{r}(e)=-e,
 \quad \half{s}(f)=-f.
\]
Apply \(\half r\) and use \cref{eq:fourcompositions}.

\begin{lemma}[Pairing with the second orbit]
\label{lem:orbitpairs}
The half-turned tile \(R(T)=\half r(T)\) has boundary paths satisfying
\begin{align}
 \half r(b)&=-\trans u(b),\label{eq:bpair}\\
 \half r(c)&=-\trans{u-t}(c),\label{eq:cpair}\\
 \half r(e)&=-e,\label{eq:epair}\\
 \half r(f)&=-\trans t(f).\label{eq:fpair}
\end{align}
Moreover the translated-pair relation \(\trans t(a)=-d\) remains a translated-pair relation after applying \(\half r\).
\end{lemma}
\begin{proof}
For \(b\), for example,
\[
 \trans u(b)=\half r\half p(b)=\half r(-b)=-\half r(b),
\]
which is \cref{eq:bpair}.  The other three calculations are identical, using the corresponding equations in \cref{eq:fourcompositions}.  Finally, applying \(\half r\) to \(\trans t(a)=-d\) and using \cref{eq:conjugate} gives
\[
 \trans{-t}(\half r(a))=-\half r(d).
\]
Thus the two paths of the turned tile pair under the lattice period \(t\), with the direction reversed as required.
\end{proof}

No angle, length, trigonometric function, or real-number limit appears.  The local matching equations are list equalities and affine polynomial identities.
